\documentclass{article}
\usepackage[T1]{fontenc}
\usepackage{lmodern}
\usepackage{graphicx}
\usepackage{amsmath,amssymb}
\usepackage[a4paper,margin=2cm]{geometry}
\usepackage[absolute,overlay]{textpos}
\setlength{\TPHorizModule}{1mm}
\setlength{\TPVertModule}{1mm}
\usepackage[indonesian]{babel}
\usepackage{hyperref}
\setlength{\emergencystretch}{2em}
% unit: Halaman 12

\begin{document}

% === Halaman 12 (llm) ===

\section*{Algoritma Knapsack Kunci-Publik}

\begin{itemize}
    \item Algoritma \textit{superincreasing knapsack} adalah algoritma yang lemah, karena cipherteks dapat didekripsi menjadi plainteksnya secara mudah dalam waktu lanjar.
    \item Algoritma \textit{non-superincreasing knapsack} atau \textit{normal knapsack} adalah kelompok algoritma \textit{knapsack} yang sulit (dari segi komputasi) karena membutuhkan waktu dalam orde eksponensial untuk memecahkannya.
    \item Namun, \textit{superincreasing knapsack} dapat dimodifikasi menjadi \textit{non-superincreasing knapsack} dengan menggunakan kunci publik (untuk enkripsi) dan kunci privat (untuk dekripsi).
\end{itemize}

\end{document}
